% L15a strategy deck: the equal-weight weighting rule.
% One short deck per rule on the strategy menu, read next to the rule's code file
% (strategies/weighting/equal_weight.jl). Written by make_strategy_decks.py (session scratchpad, 2026-11-27) and then
% kept by hand: edit this file, not the generator. Build with make in this folder: LaTeX runs from strategies/,
% so paths are relative to it.
\documentclass[aspectratio=169,11pt]{beamer}
\usepackage{vnslides}

\newcommand{\repo}{https://github.com/varnerlab/CHEME-5660-CourseRepository-Fall-2026/blob/main/lectures/week-15/L15a}

\title{\texorpdfstring{{\vnheavy L15a}}{L15a}: Equal Weight}
\subtitle{Weighting Rule, CHEME 5660 Strategy Menu, Cornell University}
\author{Copyright Jeffrey Varner 2026. All Rights Reserved}

\begin{document}

\vntitlepage

\begin{frame}{Disclaimer and Risks}
  \vspace{0.6em}
  \textbf{\ink{Training only.}} This content is for training and informational
  purposes only. The teaching team does not make, give, or endorse any offer or
  solicitation to buy or sell securities or derivative products, or provide
  investment or trading advice or strategy.

  \vspace{1.1em}
  \textbf{\ink{Risk.}} Trading involves risk and is generally not recommended for
  those with limited resources, experience, or a low risk tolerance. Past
  performance does not guarantee future success.

  \vspace{1.1em}
  \textbf{\ink{Responsibility.}} You are responsible for your investment decisions
  based on your financial circumstances, objectives, risk tolerance, and
  liquidity needs.
\end{frame}

\begin{frame}{The Idea}
  \vspace{0.3em}
  The equal-weight rule puts the same weight on every ticker in the basket:
  \begin{itemize}\setlength{\itemsep}{0.5em}
    \item \hilite{Long only}: the rule never sells a ticker short.
    \item \hilite{No estimates}: the rule ignores the preference weights, the market sensitivities, and the covariance.
    \item \hilite{Benchmark}: equal weight is a common benchmark for the other weighting rules.
  \end{itemize}

  \vspace{0.5em}
  The basket comes from the picking rule.
\end{frame}

\begin{frame}{The Equal Weights}
  \vspace{0.3em}
  For a basket of $K$ tickers, every ticker gets the same share of the wealth:
  \[
    \boxed{w_{i}=\frac{1}{K}\quad\text{for each ticker }i\text{ in the basket}}
  \]
  A basket of $K=13$ tickers gives each ticker a weight of $1/13\approx 0.077$, or 7.7\% of the wealth.
\end{frame}

\begin{frame}{Under the Trading Limits}
  \vspace{0.3em}
  At each decision, the engine scales the equal weights down until they meet every limit:
  \begin{itemize}\setlength{\itemsep}{0.5em}
    \item \hilite{Scale}: the engine finds the largest $c\leq 1$ such that the weights $c/K$ meet the limits. Every ticker keeps the same weight, and the rest of the wealth stays in cash.
    \item \hilite{Shared limits}: the cash stays at or above zero, and the weight, volatility, and trading-cost limits of every strategy apply. With the hedge on, so does the market sensitivity limit.
    \item \hilite{Hedge pool}: with the hedge on, the engine can also buy tickers with zero or negative market sensitivity from its hedge pool, as few as the limit needs.
  \end{itemize}
\end{frame}

\begin{frame}{The Scaled Weights of a Made-Up Basket}
  \vspace{0.3em}
  Here a made-up limit gives $c=0.85$ for a basket of $K=13$ tickers:
  \begin{center}
    \includegraphics{slides/figs/Fig-L15a-Deck-Weights-Equal.pdf}
  \end{center}
  The tightest limit sets the same $c$ for all 13 tickers.
\end{frame}

\begin{frame}{Reading the Code}
  \vspace{0.3em}
  The rule's code file is
  \href{\repo/strategies/weighting/equal_weight.jl}{\texttt{strategies/weighting/equal\_weight.jl}}:
  \begin{itemize}\setlength{\itemsep}{0.5em}
    \item \texttt{weigh\_equal(count)} computes the weights $1/K$ for \texttt{count} $=K$ tickers.
    \item \texttt{MyL15aEqualWeightRule} is the type the engine calls. The rule has no settings of its own.
    \item \texttt{MPC.rule\_preferences(...)} marks the picked tickers with a score of one, so that the rule can find the basket.
    \item \texttt{MPC.allocate(rule, p)} builds the equal weights of the basket, then calls \texttt{MPC.scaled\_book(...)}, which scales them to meet the limits.
  \end{itemize}

  \vspace{0.5em}
  The \href{\repo/strategies/CHEME-5660-L15a-Strategy-Menu-Fall-2026.ipynb}{strategy menu} lists every rule, and the \href{\repo/engine/CHEME-5660-L15a-Trading-Engine-Fall-2026.ipynb}{trading engine} describes the shared limits.
\end{frame}

\end{document}
